\documentclass[11pt]{article}
\usepackage[a4paper,margin=1.1in]{geometry}
\usepackage[T1]{fontenc}
\usepackage[utf8]{inputenc}
\usepackage[english]{babel}

\usepackage{lmodern}
\usepackage{amsmath,amssymb,amsthm,mathrsfs}
\usepackage{microtype}
\usepackage{fancyhdr}
\usepackage{enumitem}
\usepackage{longtable}
\usepackage{array}
\setlength{\parindent}{1.5em}
\setlength{\parskip}{0.18em}
\pagestyle{fancy}
\fancyhf{}
\cfoot{\thepage}
\newcommand{\Om}{\Omega}
\newcommand{\Omm}{\Omega_m}
\newcommand{\Hm}{H_m}
\newcommand{\Norm}{\mathrm{N}}
\newcommand{\Normone}{\mathrm{N}_1}
\newcommand{\frp}{\mathfrak p}
\newcommand{\frP}{\mathfrak P}
\newcommand{\frG}{\mathfrak G}
\newcommand{\frA}{\mathfrak A}
\newcommand{\frB}{\mathfrak B}
\newcommand{\Gg}{\mathfrak G}
\newcommand{\Aa}{\mathfrak A}
\newcommand{\Bb}{\mathfrak B}
\newcommand{\pp}{\mathfrak p}
\newcommand{\eps}{\varepsilon}
\newcommand{\srcpage}[1]{\par\medskip\noindent\textsf{[source p.~#1]}\par\smallskip}

\lhead{Heinrich Weber, Lehrbuch der Algebra, Vol. II}
\rhead{literal repair cumulative}
\begin{document}
\section*{\S\,176. The real field $H_m$ contained in $\Omega_m$.}

Every number of the field $\Omega_m$ is transformed by the substitution $(r,r^{-1})$ into the conjugate-imaginary number which is also obtained by the interchange $(i,-i)$. The product of two such conjugate-imaginary numbers is the square of the absolute value of either of them.

A number in $\Omega_m$ is real if it remains unchanged under the interchange $(r,r^{-1})$, and only under this condition. Every such number can be expressed rationally by the two-term period $r+r^{-1}$, and therefore belongs to a real field of degree $\frac12\varphi(m)$ which is a divisor of $\Omega_m$. This field we shall call the real field contained in $\Omega_m$ and denote it by $H_m$ (although other real fields are also contained in $\Omega_m$, namely all divisors of $H_m$).

The $\frac12\varphi(m)=\frac12\mu$ numbers
\begin{equation}
1,\quad r+r^{-1},\quad r^2+r^{-2},\ldots,\quad r^{\frac12\mu-1}+r^{-\frac12\mu+1}
\tag{1}
\end{equation}
form a basis of the algebraic integers of the field $H_m$. For by \S\,170, (7),
\begin{equation}
\begin{aligned}
\omega={}&x_0+x_1r+x_2r^2+\cdots+x_{\frac12\mu-1}r^{\frac12\mu-1}+x_{\frac12\mu}r^{\frac12\mu} \\
&+x_{-1}r^{-1}+x_{-2}r^{-2}+\cdots+x_{-\frac12\mu+1}r^{-\frac12\mu+1}
\end{aligned}
\tag{2}
\end{equation}
is an algebraic integer of the field $\Omega_m$ if and only if the coefficients $x_0,x_1,
\ldots$ are rational integers.

The condition that this number should belong to the field $H_m$ is obtained by carrying out the substitution $(r,r^{-1})$ and setting the difference equal to zero. If one then also divides by $r-r^{-1}$, one obtains
\begin{align*}
0={}&(x_1-x_{-1})+(x_2-x_{-2})(r+r^{-1})+\cdots \\
&+(x_{\frac12\mu-1}-x_{-\frac12\mu+1})\frac{r^{\frac12\mu-1}-r^{-\frac12\mu+1}}{r-r^{-1}}
+x_{\frac12\mu}\frac{r^{\frac12\mu}-r^{-\frac12\mu}}{r-r^{-1}}.
\end{align*}
The divisions can be carried out here; and if one then multiplies by $r^{\frac12\mu-1}$, there results for $r$ a rational equation of degree lower than the $\mu$th, which can be satisfied only if all its coefficients vanish. This leads to the equations
\[
x_{\frac12\mu}=0,\qquad x_{\frac12\mu-1}=x_{-\frac12\mu+1},\ldots,\qquad x_1=x_{-1},
\]
and hence every algebraic integer $\omega$ of the field $H_m$ is contained in the form
\[
\omega=x_0+x_1(r+r^{-1})+x_2(r^2+r^{-2})+\cdots+x_{\frac12\mu-1}(r^{\frac12\mu-1}+r^{-\frac12\mu+1}).
\]
Instead of the basis (1), if one sets
\[
\rho=r+r^{-1},
\]
one may also choose the powers of $\rho$:
\begin{equation}
1,\quad \rho,\quad \rho^2,\ldots,\quad \rho^{\frac12\mu-1}
\tag{3}
\end{equation}
as a basis of the algebraic integers of $H_m$. For the quantities (3) can be expressed integrally in terms of the quantities (1), and conversely.

The number $\rho$ is the root of an irreducible equation of degree $\frac12\mu$, which we shall denote by $\psi(\rho)=0$. If $f(r)=0$ is the equation for $r$ (\S\,169), then the identical relation
\begin{equation}
f(x)=x^{\frac12\mu}\psi\left(x+\frac1x\right)
\tag{4}
\end{equation}
holds; differentiating gives
\begin{equation}
f'(r)=r^{\frac12\mu}\psi'(\rho)(1-r^{-2}).
\tag{5}
\end{equation}
From this we can form the fundamental number $\Delta_1$ of the field $H_m$, which, since $H_m$ contains only real numbers, must be positive. If we denote by $\Normone$ the norm with respect to $H_m$, this fundamental number is
\[
\Delta_1=\pm\Normone\psi'(\rho).
\]
But if $\Norm$ is the norm formed with respect to $\Omega_m$, then
\[
\Norm\psi'(\rho)=[\Normone\psi'(\rho)]^2=\Delta_1^2,
\]
and therefore formula (5), with regard to \S\,169, gives
\[
\Delta=\Norm(1-r^{-2})\Delta_1^2.
\]
By \S\,169, (7) and (10),
\begin{equation}
\begin{aligned}
\Norm(1-r^{-2})&=q &&\text{for odd }q,\\
&=4 &&\text{for }q=2,\\
\pm\Delta&=q\Delta_1^2 &&\text{for odd }q,\\
&=4\Delta_1^2 &&\text{for }q=2.
\end{aligned}
\tag{6}
\end{equation}
If one substitutes therein
\[
\Delta=\pm q^{q^{\kappa-1}[\kappa(q-1)-1]},
\]
it follows that
\begin{equation}
\begin{aligned}
\Delta_1&=q^{\,q^{\kappa-1}\frac{q-1}{2}-\frac{q^{\kappa-1}+1}{2}} &&\text{for odd }q,\\
&=2^{(\kappa-1)2^{\kappa-2}-1} &&\text{for }q=2,
\end{aligned}
\tag{7}
\end{equation}
so that $\Delta_1$ is always a power of $q$.

\section*{\S\,177. The prime ideals in the field $H_m$.}

We now have to examine the question of the relation between the prime ideals of the field $H_m$ and those of the field $\Omega_m$.

If $\frP$ denotes any prime ideal in $H_m$ of degree $f_1$, then $\frP$ must be a divisor of a natural prime number $p$, and therefore (by \S\,171) $\frP$ can be divisible by the second or a higher power of a prime factor $\frp$ in $\Omega_m$ only if $p=q$.

If, however, $\frP$ is a divisor of $q$, then it must be a power of $\sigma=1-r$. Put, for instance,
\[
\frP=\sigma^\lambda.
\]
Taking the norm of this with respect to $\Omega_m$, which is the square of the norm with respect to $H_m$, one obtains (by \S\,169)
\[
q^{2f_1}=q^\lambda,
\]
and hence $f_1=\frac12\lambda$. Since $f_1$ is an integer, $\lambda$ must be at least equal to $2$. But $\lambda$ also cannot be greater than $2$, since $\sigma^2$ is associated with the algebraic integer $(1-r)(1-r^{-1})$ contained in $H_m$, and therefore $\sigma^2$ must be divisible by $\frP$. Consequently $f_1=1$, and we obtain the theorem:

\begin{enumerate}[label=\arabic*.,leftmargin=2.2em]
\item The prime number $q$ is the $\frac12\mu$th power of a prime number of first degree existing in $H_m$.
\end{enumerate}

Now let $p$ be different from $q$, and let $\frp$ be a prime factor of $\frP$ in the field $\Omega_m$, of degree $f$, which is carried by the substitution $(r,r^{-1})$ into $\frp'$. Then $\frP$ is divisible both by $\frp$ and by $\frp'$; on the other hand, $\frp\frp'$ is contained in $H_m$ and is therefore divisible by $\frP$. Thus we have to distinguish whether $\frp$ and $\frp'$ are distinct or not.

If $\frp$ is distinct from $\frp'$, then $\frP$ is divisible by $\frp\frp'$, and it follows that $\frP=\frp\frp'$; taking norms gives $f_1=f$.

If, however, $\frp=\frp'$, then $\frP=\frp$, and taking norms gives
\[
2f_1=f.
\]
Which of these two cases occurs depends, therefore, on whether the prime ideal $\frp$ admits the substitution $(r,r^{-1})$ or not, that is, whether $(r,r^{-1})$ is contained in the group to which the prime ideal belongs or not. By \S\,172 this difference is determined by whether there is some exponent $h$ for which the congruence
\[
p^h\equiv -1\pmod m
\]
is satisfied, or whether there is no such exponent.

We summarize the result as follows:

\begin{enumerate}[label=\arabic*.,leftmargin=2.2em,start=2]
\item The natural prime numbers $p$ different from $q$ split into two kinds, $p_1,p_2$.

The first kind $p_1$ is characterized by the condition that for some exponent $h$
\[
p_1^h\equiv -1\pmod m,
\]
and these prime numbers, if they belong to the exponent $f$ modulo $m$ and if $\mu=ef$, split in the field $H_m$ into $e$ prime factors of degree $\frac12f$. Their prime factors are not further decomposable in $\Omega_m$ either.

The prime numbers of the second kind $p_2$ are determined by the condition that none of their powers is congruent to $-1$ modulo $m$; they split in the field $H_m$ into $\frac12e$ prime factors of degree $f$. Their prime factors are still decomposable in $\Omega_m$ into two factors each.
\end{enumerate}

For prime numbers of the first kind, $f$ must certainly be even. If $m$ is odd, it is also sufficient, in order that $p$ be a prime number of the first kind, that it belong to an even exponent; for then
\[
(p^{\frac12 f}-1)(p^{\frac12 f}+1)\equiv0\pmod m.
\]
The first of these factors, $p^{\frac12 f}-1$, is not divisible by $m$, since otherwise $p$ would belong to the exponent $\frac12f$. Consequently $p^{\frac12 f}+1$ must be divisible by $q$. But here both factors $p^{\frac12 f}-1$ and $p^{\frac12 f}+1$ cannot be divisible by $q$, since otherwise their difference, which is equal to $2$, would also be divisible by $q$. Hence $p^{\frac12 f}+1$ is divisible by $m$.

If, however, $m$ is a power of $2$, then the congruence
\[
p^h\equiv -1\pmod m
\]
is possible only if $p\equiv -1\pmod m$. For every even power of an odd number is $\equiv+1\pmod 8$ and therefore cannot be $\equiv-1\pmod m$. But if $h$ is odd, then
\[
\frac{p^h+1}{p+1}=p^{h-1}-p^{h-2}+\cdots+1
\]
is itself odd, and consequently $p^h+1$ is divisible by no higher power of $2$ than $p+1$. Thus we may add to Theorem 2 the following more precise determination:

\begin{enumerate}[label=\arabic*.,leftmargin=2.2em,start=3]
\item If $m$ is odd, then the prime numbers $p_1$ belong to an even exponent, and the prime numbers $p_2$ to an odd exponent.

If $m$ is even, then the prime numbers of the form $km-1$ belong to the first kind, and all other odd prime numbers to the second kind.
\end{enumerate}

\section*{\S\,178. The units of the field $H_m$.}

The units belonging to the field $H_m$ are of special importance for the applications. The units of the field $\Omega_m$ may also be reduced to them, according to the following theorem:

\begin{enumerate}[label=\arabic*.,leftmargin=2.2em]
\item Every unit of the field $\Omega_m$ is the product of a root of unity and a real unit of the field $H_m$.
\end{enumerate}

Denote any unit of the field $\Omega_m$ by $\mathscr E(r)$. Then $\mathscr E(r^{-1})$ is the conjugate-imaginary value, which likewise represents a unit, and the quotient $\mathscr E(r):\mathscr E(r^{-1})$ is again a unit whose absolute value
\begin{equation}
\sqrt{\frac{\mathscr E(r)}{\mathscr E(r^{-1})}\frac{\mathscr E(r^{-1})}{\mathscr E(r)}}
\tag{1}
\end{equation}
is equal to $1$. Consequently this quotient is a root of unity and hence is $=\pm r^h$, where $h$ is some integral exponent (\S\,175, Theorems 1 and 2). We therefore have
\begin{equation}
\mathscr E(r)=\pm r^h\mathscr E(r^{-1}).
\tag{2}
\end{equation}

In the rest of the proof a small distinction has now to be made according as $m$ is odd or even.

First suppose that $m$ is odd. Then in (2) we may take the exponent $h$ to be even, since, if necessary, we may replace it by $h+m$. Moreover, in formula (2) only the upper sign is admissible. For if the lower sign stood there, it would follow that
\[
\mathscr E(r)\equiv\mathscr E(1)\equiv-\mathscr E(1),\qquad 2\mathscr E(1)\equiv0\,[\mathrm{mod}\,(1-r)].
\]
But by \S\,169, $1-r$ is a divisor of $q$, and so is relatively prime to $2$; therefore $\mathscr E(1)$, and consequently also $\mathscr E(r)$, is divisible by $(1-r)$. This contradicts the supposition that $\mathscr E(r)$ is to be a unit.

Consequently (2) gives
\[
r^{-\frac12 h}\mathscr E(r)=r^{\frac12 h}\mathscr E(r^{-1}),
\]
whence it follows that $r^{-\frac12 h}\mathscr E(r)$ is a real unit. Denoting it by $e(r)$, we therefore have
\begin{equation}
\mathscr E(r)=r^{\frac12 h}e(r),
\tag{3}
\end{equation}
which proves Theorem 1 in this case.

If, secondly, $m$ is a power of $2$, then $\mu=\frac12m$ and $r^\mu=-1$; and along with $r$, $-r$ is also an $m$th root of unity. If necessary replacing $h$ by $h+\frac12m$, we may assume the upper sign in (2) and write
\begin{equation}
\mathscr E(r)=r^h\mathscr E(r^{-1}).
\tag{4}
\end{equation}
The point is now to prove that $h$ must be even.

If we represent $\mathscr E(r)$ in the basis $1,r,r^2,\ldots,r^{\mu-1}$, then
\begin{equation}
\mathscr E(r)=\sum_{s=0}^{\mu-1}x_sr^s,
\tag{5}
\end{equation}
where the $x_s$ are rational integers. But by (4) we then have
\[
\sum_{s=0}^{\mu-1}x_sr^s=\sum_{s=1}^{\mu-1}x_sr^{h-s},
\]
and from this it follows that $x_s=\pm x_{s'}$ whenever $s+s'\equiv h\pmod\mu$. If $h$ is odd, then of two numbers so connected one is even and the other odd (since $\mu$ is a power of $2$). Thus the expression (5) gives for $\mathscr E(r)$ an aggregate of terms of the form
\[
x_s(r^s\pm r^{s'}).
\]
But
\[
r^s\pm r^{s'}=r^s(1\pm r^{s'-s})=r^s(1\mp r^{\mu+s'-s})
\]
is always divisible by $1-r$, and from this it would follow that $\mathscr E(r)$ is divisible by $1-r$, which is not a unit, contrary to the concept of a unit. Therefore the assumption that the exponent $h$ in formula (4) is odd is inadmissible, and, as above, we may put
\[
\mathscr E(r)=r^{\frac12h}e(r),
\]
where $e(r)$ is a real unit. Theorem 1 is thereby proved.

We shall single out a system of real units. By \S\,169 the $\mu$ quantities $1-r^n$ are all associated with one another. Therefore the quotient
\[
\frac{1-r^n}{1-r^{n'}},
\]
where $n,n'$ are two numbers relatively prime to $m$, is a unit. From this, however, setting
\[
r=e^{\frac{2\pi i}{m}},
\]
one derives the real unit
\begin{equation}
\frac{\sin\frac{n\pi}{m}}{\sin\frac{n'\pi}{m}}
=\frac{r^{\frac n2}-r^{-\frac n2}}{r^{\frac{n'}2}-r^{-\frac{n'}2}}
=r^{\frac{n'-n}{2}}\frac{1-r^n}{1-r^{n'}}.
\tag{6}
\end{equation}
If $m$ is even, one may set $n'=n+\frac m2$, and obtains in this case the real units
\begin{equation}
\tau_n=\operatorname{tang}\frac{n\pi}{m},
\tag{7}
\end{equation}
where $n$ may denote any odd number.

Replacing $n$ by $\frac12m-n$ carries $\tau_n$ into its reciprocal value. If $n$ runs through the sequence of numbers
\[
1,3,5,\ldots,\frac m4-1,
\]
then all the values $\tau_n$ obtained are positive proper fractions.

\clearpage

\srcpage{648}
\section*{Twenty-second Section. Abelian fields and cyclotomic fields.}

\section*{\S~179. Simple Abelian fields.}

We now turn to one of the most interesting applications of the theory of algebraic numbers.

The matter is the proof of the general theorem:\footnote{Kronecker first stated this theorem, but did not publish a complete proof of it. The first proof was made known by the author of the present work in vol.~8 of \textit{Acta Mathematica} (1886). Very recently a second proof by Hilbert has been published, which is based on the theorems developed in the nineteenth section (\textit{Nachrichten d. Gesellschaft d. Wissenschaften zu Göttingen}, 1896, issue 1).}

\textbf{I. All Abelian number fields in the absolute domain of rationality are cyclotomic fields.}

Once we have proved this theorem, the investigations of the third section on cyclotomic fields acquire an increased interest, because it is thereby shown that on the path indicated there one finds not only all cyclotomic fields, but all Abelian fields whatever.

Let
\[
  \Om=R(x)
\]
be an Abelian field of degree $m$, that is, a field consisting of all rational functions of a root $x$ of an irreducible Abelian equation of degree $m$, when the field $R$ of rational numbers is regarded as the domain of rationality. The Galois group $\Gg$ of this equation, which is therefore an Abelian group and has the same degree $m$ as the field $\Om$, is also the group of the field $\Om$.

If $x_1$ is a non-primitive number from $\Om$, then $\Om_1=R(x_1)$ is likewise an Abelian field, which we call a proper divisor of $\Om$, and whose degree is a proper divisor of the degree of $\Om$. For if $\Aa$ is the group to which the number $x_1$ belongs, then $\Gg\mid\Aa$ (\S~4 of this volume) is the group of the field $\Om_1$, and this is, together with $\Gg$, an Abelian group. A divisor of $\Om$ which has the same degree as $\Om$ is identical with $\Om$ (cf. Vol.~I, \S\S~144, 149, 156).

\srcpage{649}
If there are now two non-primitive numbers $x_1,x_2$ in $\Om$ of such a kind that
\[
  \Om=R(x_1,x_2)
\]
can be put, then the fields $\Om_1=R(x_1)$, $\Om_2=R(x_2)$ are proper divisors of $\Om$, and $\Om$ is called \emph{composed} from the two fields $\Om_1,\Om_2$.

If the field $\Om$ permits no such representation, it is called \emph{simple}.

Thus, if it is assumed as proved that $\Om_1,\Om_2$ are cyclotomic fields, that is, that all their numbers are representable rationally by roots of unity, then the same follows for $\Om$.

If one of the fields $\Om_1,\Om_2$ is decomposable, then we can repeat the decomposition, and, since the degrees are decreasing positive integers, we must finally arrive at simple fields.

Theorem I therefore needs to be proved only for simple Abelian fields.

It is, however, to be recalled that the decomposability of a field $\Om$ by no means already follows from the existence of a divisor, and that in a decomposable field the decomposition into simple fields can occur in quite different ways, so that for fields the laws of divisibility do not hold as they do in the domain of numbers.

A criterion for the simplicity of a field can be derived from its group.

If the group $\Gg$ has two proper divisors $\Aa$ and $\Bb$ of such a kind that every element can be composed in one and only one way from an element of $\Aa$ and an element of $\Bb$, then we call the group $\Gg$ decomposable into the two components $\Aa,\Bb$; thus $\Gg$ is decomposable when, in the product formed according to the composition of parts (\S~4),
\begin{equation}
  \Gg=\Aa\Bb,
\tag{1}
\end{equation}
each element of $\Gg$ appears once and only once. The degree of $\Gg$ is then equal to the product of the degrees of $\Aa$ and $\Bb$.

If there are no such divisors, the group $\Gg$ is called indecomposable.

\srcpage{650}
Then we can prove the theorem:

\textbf{1. If the group of an Abelian field is decomposable, then the field too is composed.}

To prove this, assume that the group $\Gg$ of the field $\Om$ is decomposable according to formula (1), and let $m,a,b$ be the degrees of $\Gg,\Aa,\Bb$. We take a number $\xi$ of the field $\Om$ belonging to the group $\Bb$, which is therefore transformed by the substitutions of $\Gg$ into $a$ different conjugate values, and likewise a $b$-valued number $\eta$ belonging to $\Aa$. We can then form a number
\[
  x=\alpha\xi+\beta\eta
\]
with rational numerical coefficients $\alpha,\beta$, which has $m$ different conjugate values, and then $\Om=R(x)$. Thus $\Om$ is composed from $R(\xi)$ and $R(\eta)$ (Vol.~I, \S~143).

If we now recall the representation of an Abelian group by a basis derived in \S\S~9, 10 of this volume, then repeated application of what has just been proved gives:

\textbf{2. Every Abelian field $\Om$ can be composed from such fields whose degree is a power of a prime number and whose group is cyclic. The degrees of these composing fields are the invariants of the group of $\Om$.}

These theorems are supplemented by the following:

\textbf{3. An Abelian field of prime-power degree with cyclic group is simple.}

If the group $\Gg$ is cyclic, then its elements can be represented by the powers of a basis, $G^\gamma$, when $\gamma$ runs through a complete system of residues modulo $m$. We obtain all divisors $\Aa$ of $\Gg$ represented by
\[
  \Aa=G^{b\gamma},
\]
when $b$ is a divisor of $m$ and $\gamma$ runs through a complete system of residues modulo $a=m:b$.

If
\[
  \Aa'=G^{b'\gamma}
\]
is a second divisor of $\Gg$ of degree $a'$ and if
\[
  b'\le b,
  \qquad a'\ge a,
\]
then, when we now assume that $m$ and hence also $a,b,a',b'$ are powers of a prime number $q$, $b'$ is a divisor of $b$, and consequently $\Aa$ is a divisor of $\Aa'$.

\srcpage{651}
If therefore $\xi$ is a number of the field $\Om$ which belongs to the group $\Aa$, then $\xi$ is a primitive number of the field $R(\xi)$ of degree $b$, and in this field every number $\xi'$ is also contained which admits the substitutions of the group $\Aa'$. Thus if $\Aa,\Bb$ and $\Aa',\Bb'$ are proper divisors of $\Gg$, then $R(\xi)$ is different from $\Om$, and $R(\xi)$ is not enlarged by the adjunction of $\xi'$. Hence $\Om$ also cannot be composed from $R(\xi)$ and $R(\xi')$, whereby 3 is proved.

\section*{\S~180. The resolvents.}

According to what has now been proved, in the further considerations, whose aim is the proof of Theorem I, we may restrict ourselves to such Abelian fields whose degree is a prime power and whose group is cyclic. For this reason it also sufficed, for the intended application, that in the preceding section we restricted ourselves to the consideration of the cyclotomic fields $\Omega_m$ in which $m$ was a prime power. We retain as much as possible of the notation used there and put
\begin{equation}
  m=q^\kappa,
\tag{1}
\end{equation}
where $q$ is a prime number and $\kappa$ a positive exponent; and in the case where $q=2$ we assume $\kappa>1$. By $n$ we denote any number not divisible by $q$. Let $r$ be a primitive $m$-th root of unity, and $\Omega_m$ the field of rational functions of $r$.

Let now $\Om=R(x)$ be a simple Abelian field of degree $m$, so that $x$ is a root of a rational irreducible cyclic equation of degree $m$. We denote the roots of this equation, in a suitable order, by
\begin{equation}
  x_0,x_1,x_2,\ldots,x_{m-1},
\tag{2}
\end{equation}
and assume $x_m=x_0$, $x_{m+1}=x_1$, and so on. Then
\[
  x_1=\Phi(x_0),\quad x_2=\Phi(x_1),\ldots,\quad x_0=\Phi(x_{m-1}),
\]
where $\Phi$ denotes an integral function with rational coefficients, and the group of the field $\Om$ consists of the powers of the substitution $(x_0,x_1)$, or of the cyclic permutations

\srcpage{652}
of the roots (2). Every cyclic function of these quantities is a rational number. What we have to prove is that the $x$ can be expressed rationally by roots of unity. If we adjoin to the field $\Om$ the $m$-th root of unity $r$, there arises a field $R(x,r)$ which contains the two fields $\Om=R(x)$ and $R(r)=\Omega_m$ as divisors.

If a number $\omega=\Phi(x,r)$ of the field $R(x,r)$ admits all substitutions $(x_0,x_1)$, $(x_0,x_2)$, \ldots, $(x_0,x_{m-1})$, then it is contained in the field $\Omega_m$; for then
\[
  m\omega=\Phi(x_0,r)+\Phi(x_1,r)+\cdots+\Phi(x_{m-1},r),
\]
and the coefficients of the powers of $r$ in this expression are not only cyclic but even symmetric functions of the roots (2), hence rational numbers. This remains true even in the case where the equation for $x$ is reduced by adjunction of $r$.

We apply this to the resolvents
\begin{equation}
  (r,x_0)=x_0+rx_1+r^2x_2+\cdots+r^{m-1}x_{m-1},
\tag{3}
\end{equation}
and, more generally, when $s$ denotes an arbitrary integral exponent,
\begin{equation}
  (r^s,x_0)=x_0+r^s x_1+r^{2s}x_2+\cdots+r^{(m-1)s}x_{m-1}.
\tag{4}
\end{equation}
By these resolvents $x_0$ itself can again be rationally expressed, namely
\begin{equation}
  m x_0=\sum_{s=0}^{m-1}(r^s,x_0),
\tag{5}
\end{equation}
and if therefore we can show that all these resolvents $(r^s,x_0)$ are cyclotomic numbers, then we have reached our goal. If we put
\begin{equation}
  (r^s,x_\lambda)=x_\lambda+r^s x_{\lambda+1}+r^{2s}x_{\lambda+2}+\cdots+r^{(m-1)s}x_{\lambda+m-1},
\tag{6}
\end{equation}
then $(r^s,x_\lambda)$ arises from $(r^s,x_0)$ by the substitution $(x_0,x_\lambda)$, and the fundamental relation holds:
\begin{equation}
  r^{\lambda s}(r^s,x_\lambda)=(r^s,x_0).
\tag{7}
\end{equation}
All these resolvents are numbers of the field $R(x,r)$; we consider preferably the following two combinations:
\begin{align}
  (r^s,x_0)^m &= \omega_s, \tag{8}\\
  (r^s,x_0)^{-1}(r,x_0)^s &= \alpha. \tag{9}
\end{align}

\srcpage{653}
These two numbers admit, as follows immediately from (7), the substitutions $(x_0,x_\lambda)$ and are therefore numbers of the field $\Omega_m$.

If $n$ is not divisible by $q$, then, because of the irreducibility of the cyclotomic equation, none of the resolvents $(r^n,x_0)$ can vanish unless they all vanish; for the equation
\[
  (r,x_0)^m=0
\]
would be an equation for $r$ with rational coefficients, admitting all substitutions $(r,r^n)$.

Likewise, none of the resolvents $(r^s,x_0)$ can vanish unless the whole family vanishes in which $s$ has a fixed greatest common divisor with $m$.

If all $(r,x_0)$ vanish, formula (9) is not applicable. Then, however, we can directly consider the resolvents $(r^q,x_0)$, which belong to a field of degree $m q^{-1}$ arising from
\[
  y_0=x_0+x_q+x_{2q}+\cdots+x_{m-q},
\]
and if these quantities too vanish, we pass to the resolvents $(r^{q^2},x_0)$, and so on.

If, however, $(r,x_0)$ does not vanish, then by formula (9) all $(r^s,x_0)$ can be represented rationally by $(r,x_0)$ and by cyclotomic numbers.

Consequently it is sufficient for our purpose if we can prove that the non-vanishing resolvents $(r,x_0)$ are cyclotomic numbers.

Let $n$ run through a complete system of incongruent numbers not divisible by $q$. Then
\[
  \sum n\equiv 0\pmod m,
\]
for the numbers $n$ split into pairs which complement one another to $m$. From (7) it follows that the product
\begin{equation}
  \prod_n (r^n,x_0)=a
\tag{10}
\end{equation}
is a number in $\Omega_m$. But since this number moreover is not changed when any of the substitutions $(r,r^n)$ is performed in it, $a$ is a rational number.

\srcpage{654}
\section*{\S~181. Preparation for the proof.}

After the developments of the preceding paragraph, for the proof of the great Theorem I it remains only to show that the number contained in $\Omega_m$,
\[
  \omega=(r,x_0)^m,
\]
is the $m$-th power of a cyclotomic number.

\textbf{1. The essential property of the numbers $\omega_n$ conjugate to $\omega$, on which the proof rests, is that}
\begin{equation}
  \omega_n^{-1}\omega_1^n=\alpha^m
\tag{1}
\end{equation}
\textbf{is the $m$-th power of a number in $\Omega_m$.}

Here $n$ may be any number not divisible by $q$. If we form, of the right- and left-hand sides of (1), the norms in the field $\Omega_m$, then, denoting by $a$ the norm of the number $\alpha$, which is a rational number, and observing that every $\omega_n$ has the same norm as $\omega$, we obtain
\[
  N(\omega)^{n-1}=a^m.
\]
If now $m$ is odd, then we can put $n=2$ in this, and obtain as a consequence of 1:

\textbf{2. The norm of $\omega$ is the $m$-th power of a rational number,}
\begin{equation}
  N(\omega)=a^m.
\tag{2}
\end{equation}
If $m$ is a power of $2$, this inference is no longer permissible. It then follows from 1 only that $N(\omega)^2$ is the $m$-th power of a rational number. Nevertheless, as is seen from (10) of the preceding paragraph, the numbers $\omega$ must, also in this case, still satisfy condition 2; but this condition is then no longer contained automatically in 1, and is added as a new property.

What we have to prove is that from 1 and 2 it follows that $\omega$ itself is the $m$-th power of a cyclotomic number, even if in a higher field.

For this it is necessary to investigate the decomposition of the numbers $\omega$ into their prime factors in the field $\Omega_m$.

The prime number $q$ is, as was proved in \S~169, associated with the $\varphi(m)$-th power of a prime number $\sigma=1-r$ existing in $\Omega_m$.

\srcpage{655}
If therefore $\sigma^k$ is the power of $\sigma$ contained in $\omega$, we put
\[
  \omega=\sigma^k\omega',
\]
where $\omega'$ is relatively prime to the prime number $q$. Formation of the norm gives, by 2,
\[
  a^m=q^k b,
\]
where $b$ is a rational number which, in lowest terms, contains the prime number $q$ neither in the numerator nor in the denominator. From this it follows that $k$ must be divisible by $m$, and hence the first result:

\textbf{3. The exponent of the prime number $\sigma$ contained in $\omega$ is always divisible by $m$.}

Let now $p$ be a prime number different from $q$, and let $\pp_\xi$ be a prime factor of $p$ in the field $\Omega_m$. Let $\xi$ run through the sequence of numbers $\xi_1,\xi_2,\ldots,\xi_e$ determined more precisely in \S~172 in the different cases, so that
\begin{equation}
  p=\pp_{\xi_1}\pp_{\xi_2}\cdots\pp_{\xi_e},
\tag{3}
\end{equation}
and assume that $\pp_\xi^{k_\xi}$ is the power of $\pp_\xi$ contained in $\omega$. The power of $\pp_{n\xi}$ contained in $\omega_n$ is then $\pp_{n\xi}^{k_\xi}$; and if we determine $n'$ from the congruence
\begin{equation}
  nn'\equiv 1\pmod m,
\tag{4}
\end{equation}
then $\pp_\xi^{k_{n'\xi}}$ is the power of $\pp_\xi$ contained in $\omega_n$. Accordingly, in
\begin{equation}
  \omega_n^{-1}\omega_1^n
\tag{5}
\end{equation}
the power $\pp_\xi^{n k_\xi-k_{n'\xi}}$ is contained, and by 1 the exponent of this power must be divisible by $m$.

Thus for every $n$ not divisible by $q$ we obtain the congruence
\begin{equation}
  n k_\xi\equiv k_{n'\xi}\pmod m,
\tag{6}
\end{equation}
in which we can now also replace $\xi$ by $n\xi$, whereby
\begin{equation}
  n k_{n\xi}\equiv k_\xi\pmod m
\tag{7}
\end{equation}
is obtained. If in this we take $\xi=1$ and put $\xi,\xi'$ for $n,n'$ and $k$ for $k_1$, then from (7) follows
\begin{equation}
  k_\xi\equiv \xi' k\pmod m,
\tag{8}
\end{equation}
where $k$ now has the same value for all $\xi$, and $\xi,\xi'$ are connected by the congruence
\[
  \xi\xi'\equiv1\pmod m.
\]

\srcpage{656}
If we take $n=p$, then $\pp_{p\xi}=\pp_\xi$ (\S~172), and so also $k_{p\xi}=k_\xi$; and from (7) and (8) follows
\begin{equation}
  k(p-1)\equiv0\pmod m.
\tag{9}
\end{equation}
Important consequences now follow from this. First, if $p-1$ is not divisible by $q$, then it follows that $k$ must be divisible by $m$, and we therefore have the theorem:

\textbf{4. The prime factors of a prime number $p$ which is not congruent to $1$ modulo $q$ are contained in $\omega$ only in such powers whose exponents are divisible by $m$.}

If, in the case where $m$ is a power of $2$, $p-1$ is divisible by $2$ but not by $4$, that is, if $p\equiv -1\pmod4$, then it follows from (9) only that $k$ is divisible by $\tfrac12m$, not that $k$ is divisible by $m$. The exponent of the power of $\pp_\xi$ contained in $\omega$ is, apart from multiples of $m$, equal to $k$; and if $k$ is divisible by $m$, everything is as in case 4. But if $k$ is divisible only by $\tfrac12m$, not by $m$, so that
\[
  k\equiv \tfrac12m\pmod m,
\]
then $k-\tfrac12m$ is divisible by $m$, and with reference to the decomposition (3) we can formulate the theorem as follows:

\textbf{5. If $m$ is a power of $2$ and $p$ is a prime number congruent to $-1$ modulo $4$, then the exponent $h$ can be so determined that all prime factors of $p$ in}
\[
  \omega\sqrt{p}^{\,-hm}
\]
\textbf{are contained only in such powers whose exponent is divisible by $m$.}

Finally let $p-1$ be divisible by $q$, or, if $m$ is even, by $4$. Let $m_1$ be the greatest common divisor of $m$ and $p-1$, and let
\[
  m=m_1m_2.
\]
Then by (9) $k$ must be divisible by $m_2$; and if we put
\[
  k=hm_2,
\]
then by (8) the product contained in $\omega$ is
\begin{equation}
  \left(\prod_\xi \pp_\xi^{\xi'}\right)^{hm_2}.
\tag{10}
\end{equation}
In addition, the prime factors $\pp_\xi$ occur in $\omega$ only in such powers whose exponents are divisible by $m$.

\srcpage{657}
In this case $\xi$, by \S~172, III, runs through a complete system of residues not divisible by $q$ modulo $m_1$, and in (10) we may understand by $\xi'$ the least positive solution of the congruence
\[
  \xi\xi'\equiv1\pmod{m_1},
\]
because, if we change the exponents $\xi'$ in (10) modulo $m_1$, only $m$-th powers of the prime factors $\pp$ are added.

Then we can apply Kummer's theorem [\S~174, (16)] to the product (10), replacing $m$ there by $m_1$, and obtain, if we understand by $r_1=r^{m_2}$ an $m_1$-th root of unity and by the $\eta$ certain periods of the $p$-th roots of unity,
\[
  \prod_\xi \pp_\xi^{\xi'}=(r_1,\eta)^{m_1},
\]
and consequently
\begin{equation}
  \left(\prod_\xi \pp_\xi^{\xi'}\right)^{hm_2}=(r_1,\eta)^{hm}.
\tag{11}
\end{equation}
Consequently we have the theorem:

\textbf{6. If $p$ is a prime number and $p-1$ is divisible by $q$, or, when $q=2$, by $4$, then the positive exponent $h$ can be so determined that the number contained in $\Omega_m$}
\[
  \omega(r_1,\eta)^{-hm}
\]
\textbf{contains the prime factors of $p$ only in such powers whose exponent is divisible by $m$.}

It is now also no longer necessary to distinguish Theorem 5 from Theorem 6; for for $m_1=2$, 6 passes into 5, because then $r_1=-1$, and by Vol.~I, \S~171,
\[
  (r_1,\eta)^m=(-1,\eta)^m=\sqrt{p^{\,m}}
\]
holds.\footnote{In the author's paper in vol.~8 of \textit{Acta Mathematica} it was omitted to emphasize separately the case of Theorem 5.}

If we collect the result of Theorems 3 to 6 in one formula, we obtain, when $\eps$ denotes a unit functional, $\varphi$ any functional of the field $\Omega_m$, and $p$ runs through a sequence of prime numbers congruent to $1$ modulo $q$, among which the same prime number $p$ may occur several times,
\begin{equation}
  \omega=(r,x_0)^m=\eps\,\varphi^m\prod_p (r^{m_2},\eta)^m.
\tag{12}
\end{equation}

\srcpage{658}
The resolvents of the cyclotomy occurring in these formulae, $(r^{m_2},\eta)$, are cyclotomic numbers in a higher field. But their $m$-th powers are numbers contained in the field $\Omega_m$ which are transformed by the substitution $(r,r^n)$ into
\[
  (r^{n m_2},\eta)^m.
\]
Likewise the combinations
\begin{equation}
  (r^{n m_2},\eta)^{-1}(r^{m_2},\eta)^n
\tag{13}
\end{equation}
are contained in the field $\Omega_m$ [\S~174, (9), (11)]. These resolvents are special cases of the general resolvent $(r,x_0)$. From (12), if we denote by $\varphi_n,\eps_n$ the functionals conjugate to $\varphi$ and $\eps$, we obtain
\begin{equation}
  \omega_n=\eps_n\varphi_n^m\prod_p (r^{m_2n},\eta)^m.
\tag{14}
\end{equation}
This formula teaches us the most important properties of the functionals $\varphi_n$ and $\eps_n$, first:

\textbf{7. The functionals $\varphi_n^m$ are associated with numbers of the field $\Omega_m$, and therefore belong to the principal class (\S~152).}

By 1, $\omega_n^{-1}\omega_1^n=\alpha^m$ is the $m$-th power of a number of $\Omega_m$. Since the combinations (13) are also numbers in $\Omega_m$, we can put
\[
  \prod_p (r^{n m_2},\eta)^{-1}(r^{m_2},\eta)^n=\frac{\alpha}{\beta},
\]
and obtain from (14)
\[
  \eps_n^{-1}\eps_1^n(\varphi_n^{-1}\varphi_1^n)^m=\beta^m,
\]
where $\beta$ is a number in $\Omega_m$.

It follows from this that the product
\[
  \beta\varphi_n\varphi_1^{-n}=\eps
\]
is a unit functional of the field $\Omega_m$, and from this arise, for the functionals $\varphi_n$ and the unit functionals $\eps_n$, the following two theorems:

\textbf{8. The conjugate functionals $\varphi_n$ have the property that all products}
\[
  \varphi_n\varphi_1^{-n}
\]
\textbf{belong to the principal class.}

\textbf{9. The unit functionals $\eps_n$ have the property that the products}
\[
  \eps_n\eps_1^{-n}=\eps^m
\]
\textbf{are $m$-th powers of unit functionals.}


\srcpage{659}
For the proof of the principal theorem I, everything now depends on the proof of the two propositions which are to be inferred from the properties of the functionals $\varphi$ and $\eps$:

\textbf{A. The functionals $\varphi_n$ belong to the principal class.}

If we may presuppose this lemma, then $\varphi_n$ is associated with a number $\alpha_n$, and we can write formula (14), with a somewhat changed meaning of $\eps_n$, in the form
\begin{equation}
  \omega_n=\eps_n\alpha_n^m\prod_p(r^{m_2n},\eta)^m.
\tag{15}
\end{equation}
In this formula, however, $\eps_n$ is a numerical unit of the field $\Omega_m$ satisfying theorem 9.

If, therefore, from theorem 9 we can still prove the second lemma:

\textbf{B. A numerical unit $\eps_n$ of the field $\Omega_m$, which satisfies theorem 9, is the $m$-th power of a unit in $\Omega_m$, multiplied by a power of $r$;}

then in (15) we can extract the $m$-th root, and obtain, when by $\rho$ a root of unity of possibly higher degree than $m$ and by $\alpha$ a number of the field $\Omega_m$ is denoted,
\begin{equation}
  (r,x_0)=\rho\alpha\prod_p(r^{m_2},\eta),
\tag{16}
\end{equation}
where now on the right there stand nothing but cyclotomic numbers, and by this Theorem I is therefore proved.
\end{document}
